% 草稿：A100 第二配置节。**cap1 落地前不要 \input 进 main.tex。**
% 数字来源：paper/data/a100.csv（gen_a100.py 生成）。
% 判据：docs/prereg_a100_addendum_20260925.md 第一节，唯一判决轴 = 兑现率穿过 50% 的 R̄。
% 纪律：判决只有一条轴；ε 是附加佐证，不许抬成判据；A100 上没跑控制器，不许暗示跑过。

\section{The Same Measurements on a Second GPU}
\label{sec:a100}

\S\ref{sec:xtrace} answers the objection that the trace family is the finding.
This section answers its hardware twin. Everything up to here is one
accelerator, and the paper's central variable---the size of the memory tier
relative to the working set---is exactly the quantity a bigger accelerator
changes. If the mechanism we name is real, it should move with the tier.

We ran the static sweep, and only the static sweep, on an A100 (80\,GB) with
the same model, the same trace (the replay file's checksum is identical on both
machines), the same engine build, and the same \texttt{mem-fraction-static}.
The pool comes out at 983{,}728 tokens, \textbf{9.70$\times$} the 3090's
101{,}432. The criterion was fixed before any of these runs finished; the
pre-registration note records the timestamps.

\subsection{What the tier buys}

\begin{table}[t]
  \centering
  \small
  \setlength{\tabcolsep}{4pt}%
  \begin{tabular}{llrrrrr}
    \toprule
    Machine & Pool & $\overline{R}_{\text{busy}}$ & Occupancy & Realized reuse & $\tau_2$ (ms) & SLO (\%) \\
    \midrule
    3090-24G & 101{,}432  & 0.9980 & 22.4\% & 92.5\% & 20.6 & 94.8 \\
    A100-80G & 971{,}456  & 0.9997 &  2.1\% & 92.5\% & 14.3 & 99.9 \\
    \midrule
    3090-24G & 101{,}432  & 1.6505 & 66.4\% & 60.7\% & 25.8 & 55.7 \\
    A100-80G & 983{,}728  & 1.9702 &  6.5\% & 92.5\% & 14.5 & 100.0 \\
    \midrule
    3090-24G & 101{,}432  & 1.8409 & 74.3\% & 36.4\% & 31.2 & 27.1 \\
    A100-80G & 983{,}728  & 2.3356 &  8.0\% & 92.5\% & 15.5 & 100.0 \\
    \bottomrule
  \end{tabular}
  \caption{The same static caps on two accelerators. Rows are paired at the
  same nominal cap: one, four and five, top to bottom. The A100's cap-one row
  came up with a pool of 971{,}456 rather than 983{,}728---see
  \S\ref{sec:a100:limits}---which does not affect its column: it uses 2.1\% of
  whatever pool it has. The realized-reuse column is
  \texttt{reuse\_tok} over the trace's own ceiling (\S\ref{sec:model}); it is
  the same quantity Table~\ref{tab:runs} is built from, computed by the same
  function.}
  \label{tab:a100}
\end{table}

At the higher of the two running counts the 3090 has lost nearly two-thirds of
its reuse and most of its SLO; the A100 is at 92.5\% reuse and a 100\%
success rate, using 8\% of its pool. Its reuse figure is not an artifact of one
row: the five A100 runs realize 19934952, 19934942, 19934944, 19935064 and
19934943 tokens respectively---close, but not identical, because they are five
separate executions of a pool that never evicts anything.

\subsection{The pre-registered test, and its answer}

The criterion was fixed in advance and has a single decision axis: the
$\overline{R}_{\text{busy}}$ at which a machine's realized reuse crosses 50\%.
On the 3090 that crossing is at $\overline{R}_{\text{busy}} = 1.7344$. The A100
does not cross: at its highest running count, 2.3356---well past the 3090's
crossing point---it is still at 92.5\%, and it is at 92.5\% at its lowest
count too, so the whole sweep is flat. The pre-registered rule therefore
returns \emph{supported}.

We report the mechanical output rather than a reading of the table, and we
note what the test does not cover in \S\ref{sec:a100:limits}.

\subsection{The step-cost elasticity is flat, not steep}

\S\ref{sec:criterion} makes the argument in terms of $\epsilon$, the
elasticity of per-step cost with respect to the running count. The same
measurement on the A100, at the same definition (\texttt{a100\_elasticity.csv};
that file's 3090 rows reproduce the device-only rows of
\texttt{elasticity\_hc.csv}):

\begin{center}
\small
\begin{tabular}{llrr}
  \toprule
  Machine & Segment & $\overline{R}_{\text{busy}}$ & $\epsilon$ \\
  \midrule
  3090-24G & cap1 $\to$ cap2         & 1.00 $\to$ 1.27 & $-0.21$ \\
  3090-24G & cap3 (fixed) $\to$ cap4 & 1.53 $\to$ 1.65 & $+1.15$ \\
  3090-24G & cap4 $\to$ cap5         & 1.65 $\to$ 1.84 & $+1.14$ \\
  \midrule
  A100-80G & cap1 $\to$ cap2         & 1.00 $\to$ 1.24 & $-0.74$ \\
  A100-80G & cap3 $\to$ cap4         & 1.60 $\to$ 1.97 & $+0.03$ \\
  A100-80G & cap4 $\to$ cap5         & 1.97 $\to$ 2.34 & $+0.34$ \\
  \bottomrule
\end{tabular}
\end{center}

The first segment is negative on both machines---$-0.21$ on the 3090 and
$-0.74$ on the A100---which is the direction the criterion expects at low
concurrency, though we measure the A100's steeper value without attributing it.
The 3090 crosses one at $\overline{R}_{\text{busy}} \approx 1.53$; the A100,
over a wider and further-right interval, does not come near it. This is the
mechanism the criterion names, moving with the tier as it should---but it is a
supporting measurement, not a second criterion, and we do not let it carry the
verdict.

\subsection{What this section does and does not buy}
\label{sec:a100:limits}

It buys the tier mechanism a second configuration. The pool grows by
$9.7\times$; realized reuse stops degrading, the step cost stops rising with
concurrency, and the engine stops being congested at all. That is the
mechanism of \S\ref{sec:criterion}---the lever an admission controller would
need---behaving as a function of the tier rather than as a fact about one
card. We stop short of the stronger claim: nothing here shows that a
controller placed on this tier would lose, because none was placed there.

One detail has to be stated rather than discovered: the cap-one run came up
with a pool of 971{,}456 tokens rather than the 983{,}728 the other four got,
because the card is shared and SGLang sizes the pool from free memory at
startup. The difference is 1.25\%; that run uses 2.1\% of whatever it has, and
its realized reuse is 0.925469---the same to six digits as the others. We
report it because the comparison rests on holding everything but the card
fixed, and that premise is very slightly false for one of the five runs.

It does \emph{not} buy a second negative result. \textbf{We did not run the
controller on the A100.} The four static caps are all that exist there, so this
section says nothing about whether \sys{} or AIMD would win on a larger
tier---that is untested, and we would need another batch of runs to say
anything about it. What the section rules out is narrower and still worth
having: the tier mechanism is not an artifact of the 24\,GB card.
